% chktex-file 1
% chktex-file 44
% chktex-file 45
% chktex-file 36
% chktex-file 3
% chktex-file 8
% chktex-file 18
% chktex-file 12

\documentclass[twoside]{article}
\usepackage{graphicx,fancyhdr,amsmath,amssymb,amsthm,subfig,url,hyperref}
\usepackage[margin=1in]{geometry}
\newtheorem{theorem}{Theorem}


%---------------- IMPORTANT: Student and Homework Information -------------

%%% PLEASE FILL THIS OUT WITH YOUR INFORMATION
\newcommand{\myname}{Muhammad Mujtaba}
\newcommand{\myid}{28100480}
\newcommand{\hwNo}{Assignment 6}
%%% END



\fancypagestyle{plain}{}
\pagestyle{fancy}
\fancyhf{}
\fancyhead[RO,LE]{\sffamily\bfseries\large LUMS}
\fancyhead[LO,RE]{\sffamily\bfseries\large CS 210 Discrete Mathematics}
\fancyfoot[LO,RE]{\sffamily\bfseries\large \myname:\myid @lums.edu.pk}
\fancyfoot[RO,LE]{\sffamily\bfseries\thepage}
\renewcommand{\headrulewidth}{1pt}
\renewcommand{\footrulewidth}{1pt}

%--------------------- This is the title of the document. DO NOT CHANGE IT ------------------------

\title{CS 210 \hwNo}
\author{\myname \qquad Student ID:\myid}
\date{}

%--------------------------------- AFTER Entering the Student and Homework Information, write your answers below  ----------------------------------

\begin{document}
\maketitle

\section*{Problem 1}
Define \( f:\mathbb{N} \to X \) by \[f(n) = -n.\]
If \( f(n_1) = f(n_2) \), then \(-n_1 = -n_2 \Rightarrow n_1 = n_2\). Hence \(f\) is injective.
For any \(x \in X\), we have \(x = -k\) for some \(k \in \mathbb{N}\). Then \(f(k) = -k = x\). Hence \(f\) is surjective.
Thus, \(f\) is a bijection and \(X\) is countably infinite.
\[\boxed{|X| = |\mathbb{N}|.}\]

\section*{Problem 2}
Let \(f:A\to B\) and \(g:B\to C\) be bijections. Consider \(h=g\circ f:A\to C\), defined by \(h(a)=g(f(a))\). \\
If \(h(a_1)=h(a_2)\), then \(g(f(a_1))=g(f(a_2))\). Since \(g\) is injective, \(f(a_1)=f(a_2)\). As \(f\) is injective, \(a_1=a_2\). \\
For any \(c\in C\), there exists \(b\in B\) with \(g(b)=c\). Since \(f\) is surjective, there exists \(a\in A\) with \(f(a)=b\). Then \(h(a)=c\).
Thus, \(h\) is bijective and \(|A|=|C|\).
\[\boxed{|A|=|C|.}\]

\section*{Problem 3}

\begin{enumerate}
    \item(i)    Sum of the first \(k\) naturals \\
                \textbf{Base case:} For \(k=1\), both sides equal 1 \\
                \textbf{Inductive hypothesis:} Assume \(1+2+\cdots+k = \tfrac{k(k+1)}{2}\). \\
                \textbf{Inductive step:} \[1+2+\cdots+k+(k+1)=\frac{k(k+1)}{2}+(k+1)=\frac{(k+1)(k+2)}{2}.\]
                Thus, by induction, the formula holds for all \(n\in\mathbb{N}\)
    \item(ii)   Sum of squares \\
                \textbf{Base case:} For \(k=1\), \(1^2=\tfrac{1\cdot2\cdot3}{6}=1\). \\
                \textbf{Inductive hypothesis:} Assume \(1^2+\cdots+k^2=\tfrac{k(k+1)(2k+1)}{6}\). \\
                \textbf{Inductive step:} \[1^2+\cdots+k^2+(k+1)^2=\frac{k(k+1)(2k+1)}{6}+(k+1)^2.\]
                \[=\frac{(k+1)(k(2k+1)+6(k+1))}{6}=\frac{(k+1)(k+2)(2k+3)}{6}.\]
                Thus, by induction, the formula holds for all \(n\in\mathbb{N}\)
\end{enumerate}

\newpage

\section*{Problem 4}
\textbf{Base case:} For \(k=1\), one line divides the plane into two regions. \[\frac{1^2+1+2}{2}=2. \Rightarrow \text{True}\]
\textbf{Inductive hypothesis:} Assume \(n_k=\tfrac{k^2+k+2}{2}\) regions for \(k\) lines. \\
\textbf{Inductive step:} Adding the \((k+1)^{th}\) line intersects all previous \(k\) lines at \(k\) distinct points, dividing the new line into \(k+1\) segments. Each segment divides an existing region into two, adding \(k+1\) new regions. Thus,
\[n_{k+1}=n_k+(k+1)=\frac{k^2+k+2}{2}+(k+1)=\frac{(k+1)^2+(k+1)+2}{2}.\] 
Thus, by induction, the formula holds for all \(n\in\mathbb{N}\).





\end{document}
